\documentclass[10pt,a4paper]{amsart}
\usepackage[margin=19mm]{geometry}
\usepackage{amsmath,amssymb,mathtools}
\usepackage[scr=rsfs,cal=euler]{mathalfa}
\usepackage{xfrac}
\usepackage[unicode,hidelinks,bookmarks=false]{hyperref}
\newcommand{\ie}{i.e.\ }
\newcommand{\cf}{cf.\ }
\newcommand{\resp}{respectively\ }
\newcommand{\Subsection}{\subsection}
\providecommand{\nobreakdash}{\nobreak\hbox{-}}
\setlength{\emergencystretch}{2em}
\multlinegap0pt
\usepackage{amssymb}
\usepackage{tikz}
\usepackage{algorithm}\usepackage{algpseudocode}
\newtheorem{lemma}{Lemma}
\newtheorem{theorem}{Theorem}
\title{Cocycle Actions on Hidden Quantum Markov Models: Symmetry Protection and Topological Order}
\author{Abdessatar Souissi, Abdessatar Barhoumi}
\date{Source: arXiv:2605.09605v2 (CC BY 4.0)}
\begin{document}
\maketitle

\noindent\emph{Source and licence.} Cocycle Actions on Hidden Quantum Markov Models: Symmetry Protection and Topological Order. Version arXiv:2605.09605v2 (CC BY 4.0), distributed by arXiv under the Creative Commons Attribution 4.0 International licence, http://creativecommons.org/licenses/by/4.0/, as recorded in the arXiv licence metadata for this exact version and shown on the abstract page https://arxiv.org/abs/2605.09605v2. Full licence notice: \texttt{licenses/arxiv-2605.09605-CC-BY-4.0.txt}.\par\smallskip

\noindent\textit{Editorial scope.} Counted unit: the article's theorem thm:global\_invariance (source0.tex lines 231--242). The article's complete original proof of it is delivered with it (source0.tex lines 244--321, one \texttt{proof} environment of 78 lines). No mathematics is written, repaired or re-derived: the statement and the proof are the article's own lines, in the article's own notation, and no retained line is substituted or re-flowed.  Retained uncounted as the source-local context the proof needs: lines 129--132; lines 134--137; lines 139--142; lines 189--195; lines 210--216.  Nothing was omitted from any retained span, so the package carries no omission record.  No retained line contains a citation, so the wrapper carries no bibliography and no bibliography entry.  No retained line contains a figure environment, an image-inclusion command or drawing code, so no image is delivered and the package declares no asset dependency.  Editorial formatting only, in addition: an AMS \texttt{amsart} preamble that restates the article's own declarations and macros used by the retained lines, namely the environments \texttt{lemma}, \texttt{theorem} on separate counters, together with the standard packages those lines need.  The retained environments print in this document as Lemma 1, Theorem 1 in order of appearance, whereas the article prints the same items with the numbering of its own full document.  The complete native source of the article as distributed in the arXiv e-print for this exact version is bundled unchanged as \texttt{sources/200238/source0.tex}.
\begin{equation}
\phi_0\bigl( \pi(g) Z \pi(g)^\dagger \bigr) = \phi_0(Z) \qquad \forall g \in G,\; \forall Z \in \mathcal{A}_H.
\label{eq:initial_invariance}
\end{equation}

\begin{equation}
\mathcal{E}_H\bigl( \pi(g) Z_1 \pi(g)^\dagger \otimes \pi(g) Z_2 \pi(g)^\dagger \bigr) = \pi(g) \, \mathcal{E}_H(Z_1 \otimes Z_2) \, \pi(g)^\dagger \qquad \forall g \in G,\; \forall Z_1, Z_2 \in \mathcal{A}_H.
\label{eq:Eh_equivariance}
\end{equation}

\begin{equation}
\mathcal{E}_{H,O}\bigl( \pi(g) X \pi(g)^\dagger \otimes \rho(g) Y \rho(g)^\dagger \bigr) = \pi(g) \, \mathcal{E}_{H,O}(X \otimes Y) \, \pi(g)^\dagger \qquad \forall g \in G,\; \forall X \in \mathcal{A}_H,\; \forall Y \in \mathcal{A}_O.
\label{eq:EOH_covariance}
\end{equation}

\begin{lemma}\label{lem:conv_cov}
Under the symmetry conditions \eqref{eq:Eh_equivariance} and \eqref{eq:EOH_covariance}, for all \(g \in G\), \(X, Z \in \mathcal{A}_H\), and \(Y \in \mathcal{A}_O\),
\begin{equation}
\mathcal{T}_{\pi(g) X \pi(g)^\dagger, \, \rho(g) Y \rho(g)^\dagger}^{(\text{conv})}\big( \pi(g) Z \pi(g)^\dagger \big) = \pi(g) \, \mathcal{T}_{X,Y}^{(\text{conv})}(Z) \, \pi(g)^\dagger.
\label{eq:conv_covariance}
\end{equation}
\end{lemma}

\begin{lemma}\label{lem:caus_cov}
Under the symmetry conditions \eqref{eq:Eh_equivariance} and \eqref{eq:EOH_covariance}, for all \(g \in G\), \(X, Z \in \mathcal{A}_H\), and \(Y \in \mathcal{A}_O\),
\begin{equation}
\mathcal{T}_{\pi(g) X \pi(g)^\dagger, \, \rho(g) Y \rho(g)^\dagger}^{(\text{caus})}\big( \pi(g) Z \pi(g)^\dagger \big) = \pi(g) \, \mathcal{T}_{X,Y}^{(\text{caus})}(Z) \, \pi(g)^\dagger.
\label{eq:caus_covariance}
\end{equation}
\end{lemma}

\begin{theorem}\label{thm:global_invariance}
Let \(\Xi = (\phi_0, \mathcal{E}_H, \mathcal{E}_{H,O})\) be an HQMM satisfying the symmetry conditions \eqref{eq:initial_invariance}, \eqref{eq:Eh_equivariance}, and \eqref{eq:EOH_covariance}. Let \(\varphi^{(\text{conv})}\) and \(\varphi^{(\text{caus})}\) be the global states on \(\mathcal{A}_{\mathbb{N}_0}\) generated by the conventional and causal transition maps, respectively, via the construction of Theorem 1. Then for all \(g \in G\), \(A \in \mathcal{A}_{H,\mathbb{N}_0}\), and \(B \in \mathcal{A}_{O,\mathbb{N}_0}\),
\begin{equation}
\varphi^{(\text{conv})}\big( \alpha_g^{(H)}(A) \otimes \beta_g(B) \big) = \varphi^{(\text{conv})}(A \otimes B),
\label{eq:global_conv_invariance}
\end{equation}
\begin{equation}
\varphi^{(\text{caus})}\big( \alpha_g^{(H)}(A) \otimes \beta_g(B) \big) = \varphi^{(\text{caus})}(A \otimes B).
\label{eq:global_caus_invariance}
\end{equation}
In particular, both global states are invariant under the combined projective-linear action of \(G\). Consequently, the hidden marginals \(\varphi_H^{(\text{conv})}\) and \(\varphi_H^{(\text{caus})}\) are invariant under the projective action \(\alpha^{(H)}\), and the observable marginals \(\varphi_O^{(\text{conv})}\) and \(\varphi_O^{(\text{caus})}\) are invariant under the linear action \(\beta\).
\end{theorem}

\begin{proof}
We present the proof for the conventional structure; the proof for the causal structure is entirely analogous, with Lemma \ref{lem:caus_cov} replacing Lemma \ref{lem:conv_cov}.

Let \(A = \bigotimes_{k=0}^n X_k \in \mathcal{A}_{H,\mathbb{N}_0}\) and \(B = \bigotimes_{k=0}^n Y_k \in \mathcal{A}_{O,\mathbb{N}_0}\) be elementary tensors supported on sites \(0, \dots, n\). The extension to arbitrary local observables follows by linearity and continuity, and the inductive limit construction ensures the result holds for all quasi-local observables.

Recall from Theorem 1 that the finite-volume state on \(\mathcal{A}_{[0,n]}\) is given by
\begin{equation}
\varphi_n^{(\text{conv})}\!\left( \bigotimes_{k=0}^n (X_k \otimes Y_k) \right) = \phi_0 \circ \mathcal{T}_{X_0,Y_0}^{(\text{conv})} \circ \mathcal{T}_{X_1,Y_1}^{(\text{conv})} \circ \cdots \circ \mathcal{T}_{X_n,Y_n}^{(\text{conv})}(\mathbb{I}_{\mathcal{A}_H}),
\label{eq:finite_volume_conv}
\end{equation}
where the composition denotes successive application of the sliced maps. In what follows, for $(\text{c})\in \{(\text{conv}), (\text{caus}) \}$  we denote the composition of the transition  maps $\mathcal{T}^{(\text{c})} _{X_j,Y_j}$ by

$$
\bigcirc_{k=m}^n \mathcal{T}_{X_k, Y_k} := \mathcal{T}_{X_m,Y_m}^{(\text{c})} \circ \mathcal{T}_{X_1,Y_1}^{(\text{c})} \circ \cdots \circ \mathcal{T}_{X_n,Y_n}^{(\text{c})}
$$

We prove by induction on \(n\) that for all \(g \in G\),
\begin{equation}
\bigcirc_{k=0}^n \mathcal{T}_{\pi(g) X_k \pi(g)^\dagger, \, \rho(g) Y_k \rho(g)^\dagger}^{(\text{conv})}(\mathbb{I}_{\mathcal{A}_H})
= \pi(g) \left( \bigcirc_{k=0}^n \mathcal{T}_{X_k,Y_k}^{(\text{conv})}(\mathbb{I}_{\mathcal{A}_H}) \right) \pi(g)^\dagger.
\label{eq:induction_claim}
\end{equation}


For a single site, Lemma \ref{lem:conv_cov} with \(Z = \mathbb{I}_{\mathcal{A}_H}\) together with the observation that \(\pi(g) \mathbb{I}_{\mathcal{A}_H} \pi(g)^\dagger = \mathbb{I}_{\mathcal{A}_H}\) yields
\[
\mathcal{T}_{\pi(g) X_0 \pi(g)^\dagger, \, \rho(g) Y_0 \rho(g)^\dagger}^{(\text{conv})}(\mathbb{I}_{\mathcal{A}_H})
= \pi(g) \, \mathcal{T}_{X_0,Y_0}^{(\text{conv})}(\mathbb{I}_{\mathcal{A}_H}) \, \pi(g)^\dagger.
\]


Assume the claim holds for chains of length \(n\). For length \(n+1\), define the composition for the last \(n+1\) sites:
\[
\mathcal{C}_n := \bigcirc_{k=1}^{n+1} \mathcal{T}_{X_k,Y_k}^{(\text{conv})}, \qquad
\mathcal{D}_n := \bigcirc_{k=1}^{n+1} \mathcal{T}_{\pi(g) X_k \pi(g)^\dagger, \, \rho(g) Y_k \rho(g)^\dagger}^{(\text{conv})}.
\]
By the induction hypothesis applied to the sites \(1, \dots, n+1\),
\begin{equation}
\mathcal{D}_n(\mathbb{I}_{\mathcal{A}_H}) = \pi(g) \, \mathcal{C}_n(\mathbb{I}_{\mathcal{A}_H}) \, \pi(g)^\dagger.
\label{eq:induction_hyp}
\end{equation}

Now compute the full composition for \(n+1\) sites:
\[
\begin{aligned}
&\bigcirc_{k=0}^{n+1} \mathcal{T}_{\pi(g) X_k \pi(g)^\dagger, \, \rho(g) Y_k \rho(g)^\dagger}^{(\text{conv})}(\mathbb{I}_{\mathcal{A}_H}) \\
&= \mathcal{T}_{\pi(g) X_0 \pi(g)^\dagger, \, \rho(g) Y_0 \rho(g)^\dagger}^{(\text{conv})}\Big( \mathcal{D}_n(\mathbb{I}_{\mathcal{A}_H}) \Big) \\
&= \mathcal{T}_{\pi(g) X_0 \pi(g)^\dagger, \, \rho(g) Y_0 \rho(g)^\dagger}^{(\text{conv})}\Big( \pi(g) \, \mathcal{C}_n(\mathbb{I}_{\mathcal{A}_H}) \, \pi(g)^\dagger \Big) \quad \text{by \eqref{eq:induction_hyp}} \\
&= \pi(g) \, \mathcal{T}_{X_0,Y_0}^{(\text{conv})}\big( \mathcal{C}_n(\mathbb{I}_{\mathcal{A}_H}) \big) \, \pi(g)^\dagger \quad \text{by Lemma \ref{lem:conv_cov}} \\
&= \pi(g) \left( \bigcirc_{k=0}^{n+1} \mathcal{T}_{X_k,Y_k}^{(\text{conv})}(\mathbb{I}_{\mathcal{A}_H}) \right) \pi(g)^\dagger.
\end{aligned}
\]
This completes the induction.

Now apply the initial state \(\phi_0\). By the invariance condition \eqref{eq:initial_invariance},
\[
\begin{aligned}
\varphi_n^{(\text{conv})}\big( \alpha_g^{(H)}(A) \otimes \beta_g(B) \big)
&= \phi_0 \left( \bigcirc_{k=0}^n \mathcal{T}_{\pi(g) X_k \pi(g)^\dagger, \, \rho(g) Y_k \rho(g)^\dagger}^{(\text{conv})}(\mathbb{I}_{\mathcal{A}_H}) \right) \\
&= \phi_0 \left( \pi(g) \left( \bigcirc_{k=0}^n \mathcal{T}_{X_k,Y_k}^{(\text{conv})}(\mathbb{I}_{\mathcal{A}_H}) \right) \pi(g)^\dagger \right) \\
&= \phi_0 \left( \bigcirc_{k=0}^n \mathcal{T}_{X_k,Y_k}^{(\text{conv})}(\mathbb{I}_{\mathcal{A}_H}) \right) \\
&= \varphi_n^{(\text{conv})}(A \otimes B).
\end{aligned}
\]

Since this holds for all finite \(n\) and the embeddings \(\iota_{m,n}\) are compatible with the group actions (i.e., \(\iota_{m,n} \circ (\alpha_g^{(H)} \otimes \beta_g) = (\alpha_g^{(H)} \otimes \beta_g) \circ \iota_{m,n}\)), the projective limit state \(\varphi^{(\text{conv})}\) on \(\mathcal{A}_{\mathbb{N}_0}\) satisfies \eqref{eq:global_conv_invariance}.

\noindent\textit{Proof for the causal structure.} The proof follows exactly the same inductive structure, with \(\mathcal{T}^{(\text{caus})}\) replacing \(\mathcal{T}^{(\text{conv})}\) and Lemma \ref{lem:caus_cov} replacing Lemma \ref{lem:conv_cov}. The finite-volume state is
\[
\varphi_n^{(\text{caus})}\!\left( \bigotimes_{k=0}^n (X_k \otimes Y_k) \right) = \phi_0 \circ \mathcal{T}_{X_0,Y_0}^{(\text{caus})} \circ \mathcal{T}_{X_1,Y_1}^{(\text{caus})} \circ \cdots \circ \mathcal{T}_{X_n,Y_n}^{(\text{caus})}(\mathbb{I}_{\mathcal{A}_H}).
\]
One proves by induction that
\[
\bigcirc_{k=0}^n \mathcal{T}_{\pi(g) X_k \pi(g)^\dagger, \, \rho(g) Y_k \rho(g)^\dagger}^{(\text{caus})}(\mathbb{I}_{\mathcal{A}_H})
= \pi(g) \left( \bigcirc_{k=0}^n \mathcal{T}_{X_k,Y_k}^{(\text{caus})}(\mathbb{I}_{\mathcal{A}_H}) \right) \pi(g)^\dagger,
\]
using Lemma \ref{lem:caus_cov} at each step. Applying \(\phi_0\) and its invariance \eqref{eq:initial_invariance} yields \eqref{eq:global_caus_invariance}.
\end{proof}

\end{document}
